% PS3 mathematical companion. Build with make in this directory.
% The study order follows the assignment README; no reference results appear here.
\documentclass[aspectratio=169,11pt]{beamer}
\usepackage{vnslides}
\newcommand{\bw}{\mathbf{w}}
\newcommand{\bone}{\mathbf{1}}
\newcommand{\Sg}{\widehat{\boldsymbol{\Sigma}}_g}
\newcommand{\ghat}{\widehat{\mathbf{g}}}
\newcommand{\Var}{\operatorname{Var}}
\newcommand{\Cov}{\operatorname{Cov}}
\title{\texorpdfstring{{\vnheavy PS3}}{PS3}: Estimation Windows and Portfolio Risk}
\subtitle{CHEME 4/5660 Cornell University | Mathematical Companion}
\author{Copyright Jeffrey Varner 2026. All Rights Reserved}
\begin{document}
\vntitlepage

\begin{frame}{Disclaimer and Risks}
\textbf{Training only.} This assignment uses historical prices and simplified
investment assumptions for teaching. It provides no investment advice.

\vspace{1em}
\textbf{Short positions.} Borrowed shares must be returned. A rising stock
price can produce losses larger than the initial short-sale proceeds.

\vspace{1em}
\textbf{Exercise assumptions.} Fractional shares are allowed. Short-sale
proceeds finance long holdings. We model borrowing fees, while omitting
dividends, other trading costs, collateral requirements, and forced liquidation.
\end{frame}

\begin{frame}{The Portfolio Question}
Does using more recent data lead to a better portfolio, and does allowing
short positions help?
\vspace{0.7em}
\begin{itemize}\setlength{\itemsep}{0.7em}
\item \hilite{Choose the allocations.} Compare estimated risk with and without
short positions, using 2014--2025 and then 2025 alone.
\item \hilite{Test the holdings.} Evaluate their 2026 wealth and risk, then
calculate what borrowing the shares costs.
\item \hilite{Assess the evidence.} Interpret the observed comparison.
Advanced adds modeled probabilities of beating the benchmark.
\end{itemize}
\slidenote{Reading guide}{Studies 1--3 are shared. The final MAGBM study is Advanced only.}
\end{frame}

\begin{frame}{Study 1: Change the Estimation Window}
Keep the eight stocks, purchase date, and optimization rules unchanged.
\begin{itemize}\setlength{\itemsep}{0.6em}
\item \hilite{2014--2025 window:} all 3,017 prices from 2014--2025, giving
3,016 consecutive growth observations.
\item \hilite{2025 window:} first select the 250 prices dated in 2025,
then calculate 249 growth observations.
\end{itemize}
\vspace{0.6em}
The 2025 window excludes the change from December 31, 2024 to January 2,
2025. The 2014--2025 window includes it. Both end at the same purchase prices.
\slidenote{Comparison}{Choose long-only and shorts-allowed GMV for each window. Equal weight supplies one historical reference.}
\end{frame}

\begin{frame}{Study 1: Estimate the Growth Rates}
Within the selected window, let $S_{i,k}$ be stock $i$'s price in USD/share
at observation $k$, for $i=1,\ldots,8$. With $\Delta t=1/252$ trading year, its growth rate is:
\[
 g_{i,k}=\frac{\log(S_{i,k}/S_{i,k-1})}{\Delta t},
 \qquad k=1,\ldots,n.
\]
Use $n=3016$ for 2014--2025 or $n=249$ for 2025 alone. Write
$\mathbf g_k$ for the growth vector. The sample mean and covariance are:
\[
 \ghat=\frac1n\sum_{k=1}^n\mathbf g_k,
 \qquad
 \Sg=\frac1{n-1}\sum_{k=1}^n
 (\mathbf g_k-\ghat)(\mathbf g_k-\ghat)^\top.
\]
\slidenote{Units}{Growth is in inverse years; growth-rate covariance is in inverse years squared.}
\end{frame}

\begin{frame}{Study 1: Choose the Allocations}
Let $w_i$ be stock $i$'s fraction of initial wealth, and let $\bone$ be the
vector of eight ones. Both optimized portfolios solve:
\[
 \min_{\bw}\;\bw^\top\Sg\bw,
 \qquad \bone^\top\bw=1.
\]
\begin{itemize}\setlength{\itemsep}{0.5em}
\item \hilite{Long only:} also require $w_i\geq0$ for every stock.
\item \hilite{Shorts allowed:} permit negative weights, with no position bounds.
\item \hilite{Equal weight:} use $w_i=1/8$ as the reference allocation.
\end{itemize}
\vspace{0.2em}
Estimated growth $\hat g_p=\ghat^\top\bw$ and risk $\widehat\sigma_p=\sqrt{\bw^\top\Sg\bw}$,
both in inverse years, use the lecture's linear growth approximation at the
initial weights.
\slidenote{Scope}{Add no growth target or borrowing fee. The supplied long-only floor $R=\min_i\hat g_i$ excludes no allocation.}
\end{frame}

\begin{frame}{Study 1: Solve the Shorts-Allowed Problem}
For a positive-definite sample covariance $\Sg$ and the vector of ones
$\bone$, the global minimum-variance allocation is:
\[
 \Sg\mathbf x=\bone,
 \qquad
 \bw_{\mathrm{shorts}}=\frac{\mathbf x}{\bone^\top\mathbf x}.
\]
Solve the linear system for $\mathbf x$; do not explicitly form the inverse.
In Julia, use the backslash operator with the starter's \texttt{one\_vector},
which has one entry per stock.

\vspace{0.8em}
Every long-only allocation is feasible when shorts are allowed. The
shorts-allowed minimum estimated variance therefore cannot be larger.

\vspace{0.8em}
That conclusion concerns the \hilite{same estimated covariance}. It does
not guarantee lower realized risk on the 2026 prices.
\end{frame}

\begin{frame}{Study 1: Why Short a High-Risk Stock?}
For fixed initial weights, define the weighted growth of the long positions
as $G_+$ and the growth of the shares owed as $G_-$:
\[
 G_+=\sum_{i:w_i>0}w_i g_i,
 \qquad G_-=\sum_{i:w_i<0}|w_i|g_i.
\]
The portfolio's linear growth approximation is $G_+-G_-$. Its variance is:
\[
 \Var(G_+-G_-)=\Var(G_+)+\Var(G_-)
              -\underbrace{2\Cov(G_+,G_-)}_{\text{offsetting movement}}.
\]
Positive covariance provides an offset, while the short positions also
contribute their own variance. \hilite{A stock's own risk alone does
not establish whether a short is a useful hedge.}
\end{frame}

\begin{frame}{Study 1: Interpret the Financing}
For weights $w_i$ summing to one, define long exposure $L$, short exposure
$H$, and gross exposure $E$, all relative to initial wealth:
\[
 L=\sum_i\max(w_i,0),\qquad
 H=\sum_i\max(-w_i,0),\qquad E=L+H.
\]
The budget constraint gives $L-H=1$. A short finances additional long holdings.

\vspace{0.8em}
For example, with initial wealth of USD 10,000, $L=1.20$ and $H=0.20$ mean:
\begin{itemize}\setlength{\itemsep}{0.35em}
\item USD 12,000 of long holdings;
\item shares owed initially worth USD 2,000;
\item USD 14,000 of gross exposure and USD 10,000 of net wealth.
\end{itemize}
\end{frame}

\begin{frame}{Study 2: Hold the Shares}
Let $W_0=10{,}000$ USD and $S_{i,0}$ be the December 31, 2025 price.
The signed share count $q_i$ and wealth before fees are:
\[
 q_i=\frac{W_0w_i}{S_{i,0}},
 \qquad
 W_k^{\mathrm{gross}}=\sum_i q_i S_{i,k}.
\]
A negative $q_i$ represents shares owed. \hilite{Share counts remain fixed}
through the closing day, day $N=126$. Weights change with prices.

\vspace{0.8em}
Day 0 is the final 2025 observation. Day 1 is the first 2026 observation.
An $N$-day evaluation uses $N+1$ price rows, including day 0.
\slidenote{Terminology}{Gross wealth means wealth before borrowing fees. Gross exposure measures the magnitude of all positions.}
\end{frame}

\begin{frame}{Study 2: Measure Realized Risk}
For positive gross wealth at every observation, define the portfolio's
realized growth $g_{p,k}$ and its average $\bar g_p$ over $N$ intervals:
\[
 g_{p,k}=\frac{\log(W_k^{\mathrm{gross}}/W_{k-1}^{\mathrm{gross}})}{\Delta t},
 \qquad
 \bar g_p=\frac1N\sum_{k=1}^N g_{p,k}.
\]
The realized growth-rate standard deviation is:
\[
 \sigma_{p,\mathrm{realized}}
 =\sqrt{\frac1{N-1}\sum_{k=1}^N(g_{p,k}-\bar g_p)^2}.
\]
Report it in inverse years, using wealth \hilite{before fees}. If gross
wealth is ever nonpositive, this log-growth measure is unavailable.
\end{frame}

\begin{frame}{Study 2: Separate Profit from Hedging}
For signed share count $q_i$ and closing day $N$, stock $i$'s dollar gain or
loss, $\Pi_i$, is:
\[
 \Pi_i=q_i(S_{i,N}-S_{i,0}),
 \qquad
 \Pi_{\mathrm{long}}=\sum_{i:q_i>0}\Pi_i,
 \qquad
 \Pi_{\mathrm{short}}=\sum_{i:q_i<0}\Pi_i.
\]
The two sides add to the gain before fees:
$\Pi_{\mathrm{long}}+\Pi_{\mathrm{short}}=W_N^{\mathrm{gross}}-W_0$.
A short gains when its price falls and loses when it rises. If the price
doubles, the short loses its entire initial value.

\vspace{0.8em}
Compare \hilite{terminal wealth} and \hilite{realized risk} separately. A
short can lose money while offsetting fluctuations in the long holdings.

\vspace{0.8em}
The two optimized portfolios also have different long positions. Their
performance difference cannot be attributed entirely to the shorts.
\end{frame}

\begin{frame}{Study 2: Accumulate Borrowing Fees}
Let $B_k$ be the market value of the shares owed at the start of interval
$k$, in USD. For annual borrowing rate $b$, total fees through day $N$ are:
\[
 B_k=\sum_{i:q_i<0}|q_i|S_{i,k},
 \qquad
 C_N(b)=b\Delta t\sum_{k=0}^{N-1}B_k.
\]
Use $b=0$ and $b=0.03$ with the same holdings and price path.
\begin{itemize}\setlength{\itemsep}{0.5em}
\item Charge on beginning-of-interval values; exclude the closing-day value.
\item Accrue fees without interest and pay them on the closing day.
\item Portfolios without short positions owe zero fees.
\end{itemize}
\slidenote{Convention}{These are illustrative annual rates using 252 trading days per year.}
\end{frame}

\begin{frame}{Study 2: Net Wealth and Benchmark Success}
Net terminal wealth is gross wealth less borrowing fees. For benchmark
rate $g_y=0.05$ per trading year and holding time $T=N\Delta t$, define:
\[
 W_N^{\mathrm{net}}=W_N^{\mathrm{gross}}-C_N(b),
 \qquad W_N^{\mathrm{benchmark}}=W_0 e^{g_yT}.
\]
Scaled net present value $\rho_N$ discounts the proceeds to day 0 and
compares them with the initial investment:
\[
 \rho_N=\underbrace{\frac{W_N^{\mathrm{net}}}{W_0}}_{\text{wealth ratio}}
       \underbrace{e^{-g_yT}}_{\text{discount factor}}-1.
\]
Success means $\rho_N>0$. Equality does not count. Increasing the fee
reduces net wealth while leaving the holdings and risk before fees unchanged.
\end{frame}

\begin{frame}{Study 3: Interpret the Window Comparison}
Each window estimates its own covariance. Its GMV allocation minimizes
variance under \hilite{that estimate}. Mean growth does not affect these
weights because no growth target binds.

\vspace{0.8em}
Compare all four allocations over the same 126 trading days of 2026.
Use realized risk and net wealth to assess what carried over.

\vspace{0.8em}
Changing the window changes both sample size and recency. One shared market
episode cannot establish which window or shorting rule will always work best.
\end{frame}

\begin{frame}{Advanced: Construct the MAGBM Model}
Fit one model to each window. From its $\ghat$ and $\Sg$, define the covariance
rate $\widehat C$ (\texttt{Ĉ} in the starter code), its lower Cholesky factor $A$,
and the price drift $\mu_i$:
\[
 \widehat C=\Delta t\Sg,\qquad AA^\top=\widehat C,\qquad
 \mu_i=\hat g_i+\tfrac12\widehat C_{ii}.
\]
With independent standard Brownian motions $W_j(t)$, stock $i$ follows:
\[
 \frac{dS_i(t)}{S_i(t)}=\mu_i\,dt+\sum_{j=1}^8 A_{ij}\,dW_j(t).
\]
$\widehat C$ has units of inverse years; $A$ has units of inverse square-root years.
Keep the benchmark out of the fitted price drifts.
\end{frame}

\begin{frame}{Advanced: Simulate Joint Price Paths}
The supplied simulator passes your model to the course sampler. At interval
$k$, it draws an eight-component standard normal vector $\mathbf z_k$ with
independent components, zero mean, and identity covariance. The exact daily
MAGBM update is:
\[
 \log\!\left(\frac{\mathbf S_{k+1}}{\mathbf S_k}\right)
 =\ghat\Delta t+A\sqrt{\Delta t}\,\mathbf z_k.
\]
Ratios and logarithms are componentwise. Draws are independent across days
and paths. The matrix $A$ introduces the fitted cross-stock dependence.

\vspace{0.8em}
The simulator generates 20,000 paths through day 126 per model. Within each
model, all portfolios and fee rates share the price paths. Across models, it
reuses the normal draws; each fitted model transforms them differently.
\slidenote{Supplied}{You write no simulation code. The forecasts use historical fitted drifts, not risk-neutral prices.}
\end{frame}

\begin{frame}{Advanced: Estimate the Success Probability}
Let $s$ index the $K=20{,}000$ simulated paths. On each path, the report
values your fixed shares and subtracts that path's borrowing fees. Your
function then estimates:
\[
 \hat p_N(b)=\frac1K\sum_{s=1}^K
 \mathbf1\!\left[W_{N,s}^{\mathrm{net}}(b)>W_0e^{g_yN\Delta t}\right].
\]
The indicator is one for success and zero otherwise. Keep every path in
the denominator, including nonpositive wealth outcomes.

\vspace{0.8em}
A signed sum of lognormal asset prices is generally not lognormal, so the
report computes wealth from the joint asset paths. With no forced liquidation,
it keeps the holdings through the closing day even if wealth crosses zero earlier.
\end{frame}

\begin{frame}{Advanced: Interpret the Probability Estimate}
For estimated success probability $\hat p_N$ based on $K$ independent paths,
the Monte Carlo standard error is:
\[
 \operatorname{SE}(\hat p_N)=\sqrt{\frac{\hat p_N(1-\hat p_N)}{K}}.
\]
Multiply by 100 to report it in percentage points.

\vspace{0.7em}
A small standard error means the simulation estimates the \hilite{fitted
model's probability} precisely. It does not establish that the model or
its estimated parameters describe future prices correctly.

\vspace{0.7em}
Using common paths, a higher borrowing rate cannot raise the estimated
success probability. It may leave the estimate unchanged if no outcome
crosses the benchmark threshold.
\end{frame}

\begin{frame}{Summary}
The studies compare two estimation windows on one common 2026 episode,
with and without short positions and borrowing fees.

\vspace{0.6em}
\textbf{Key takeaways}
\begin{itemize}\setlength{\itemsep}{0.55em}
\item Covariance determines whether short positions help offset portfolio
fluctuations; a stock's own risk alone is insufficient.
\item Lower estimated variance does not guarantee lower realized risk or
more terminal wealth. Borrowing fees reduce the proceeds.
\item MAGBM supplies a model-based success probability. Simulation precision
and forecast accuracy are different questions.
\end{itemize}
\vspace{0.6em}
Use these distinctions to answer the central portfolio question in PS3.
\end{frame}
\end{document}
